\section{Implementation sequence and acceptance gates}
\label{sec:development}

\subsection{First ship a probability bridge, not four new tactic names}

Forge's existing recommendation to obtain one narrow end-to-end checked lane
remains sensible~\cite{forge-conclusion}. The proposed additions should not
interrupt that baseline work. Within this extension, the smallest useful first
milestone is a \emph{positive finite transport certificate connected to an
actual source distribution}. It needs finite sums, nonnegativity, and an
explicit source equality; it does not need infinite plays, a greatest fixed
point, or probability measures on infinite paths.

The first source theorem should be the two-coins/binomial relation from
Section~\ref{sec:transport}, written using genuine probability distributions.
The input marginals must be derived from the source expressions by checked
\code{pure}/\code{bind} equations. The accepted theorem must mention those
original expressions. A theorem saying merely that a literal rational matrix
has matching row sums does not close this milestone.

\subsection{Six concrete delivery gates}

\paragraph{Gate 1: finite rational reflection.}
Implement validated tables, exact sum checking, and a theorem that accepted
tables denote probability distributions. Establish an equivalence between the
source finite type and its index set. Test zero mass, duplicate indices,
incorrect normalization, and a source enumeration missing one constructor.
Acceptance requires both a successful source theorem and rejection of the
same certificate against an altered marginal.

\paragraph{Gate 2: transport and composition.}
Prove the full-marginal checker sound, then the Hall lower bound, then the
finite \code{bind} composition rule. Check a positive example, a positive-defect
obstruction, and a two-step composition whose second coupling depends on the
first state pair. The last test prevents an implementation from supporting
only constant transition kernels while claiming compositional reasoning.

\paragraph{Gate 3: simulation witnesses.}
Implement source-owned dependent action tables and replay of the ordered
deletion trace. Prove greatestness, not only post-fixedness. The acceptance
corpus must include the $\forall a\exists b$ example, a removed initial pair,
an empty base relation, and a deliberately incomplete target-action obstruction
list. A simulation exclusion must not be rendered as a source trace-inequality
proof.

\paragraph{Gate 4: expected duration.}
Prove the finite-horizon Bellman bound, the tight-policy result, and the
path/trap result. Connect the finite-horizon definition to the chosen source
expectation semantics. Required source tests include the two-action retry
example, an unreachable trap, a reachable trap, an initially terminal state,
and a finite exact expectation that exceeds a proposed user budget. Checking
only stationary schedulers is insufficient for the stated universal theorem.

\paragraph{Gate 5: parity.}
Prove strategy closure and the threshold-rank theorem for infinite sequences.
Transfer an explicit finite source arena to its table. Check an infinite even
cycle, a favorable priority appearing only once, two adversarial strategies,
and a source-state renaming. Temporal formula translation is a separate gate;
accepting an automaton as input is not an LTL synthesis implementation.

\paragraph{Gate 6: scale and composition.}
Add untrusted Oink and Storm adapters, sparse encodings, cancellation, and
incremental search only after the semantic gates pass. Compare the adapters
on identical reified source problems and identical checker versions. An
adapter must be removable without changing the meaning of retained proofs.
Add a composite example where a relational simulation supplies a state
invariant used by a Bellman worker, or a synthesized finite strategy changes
the nondeterministic model subsequently analyzed for expected cost.

\subsection{A useful composite target}

Consider an implementation that retries a randomized operation using one of
several equivalent sampling routines. The specification uses a simpler sampler.
A transport certificate can establish each routine's distributional relation;
a simulation certificate can justify matching every implementation choice by
a specification transition; and a Bellman potential on the resulting source
model can establish a uniform expected number of retries.

The composition must record exactly which property transfers. Simulation by
itself need not preserve a particular cost unless the base relation and the
transition theorem include that cost. For expected duration, a relation that
matches one implementation step with one specification step and preserves the
terminal predicate is a suitable initial target. Stuttering or asynchronous
simulation would require an additional accounting potential. This is a concrete
place to combine the workers without silently converting one theorem into
another.

\subsection{Evaluation after Lean integration}

The first empirical comparison should not ask whether \grind{} is a parity
solver. It should ask whether adding a worker closes natural source-level goals
that the same Forge controller, premise set, and local tactics leave open.
Hold constant the source bridge, retrieved lemmas, proof budget, and definitions.
For each family, compare the full worker with an intentionally restricted
alternative on the same goals: independent-product couplings, non-alternating
action matching, or pointwise instead of expected drift. Record failure as
\status{UNKNOWN} unless the restricted method has a checked obstruction for
its own exact fragment.

Measure search time, source reification time, certificate size, elaboration
time, kernel checking time, and ordinary recompilation separately. Include
false goals, malformed certificates, and equivalent source programs with
renamed constructors. Keep theorem families disjoint between tuning and
measurement. A tiny generated model accepted by a Python checker is not a
replacement for any row of this source-level evaluation.

\subsection{Research directions after the implemented fragment}

Three expansions have especially clear certificate paths. First, min-cost
couplings and dual potentials can replace Boolean relations by quantitative
distances without changing the marginal checker. Second, symbolic drift can
reuse Forge's existing exact polynomial machinery after a proved expected
substitution. Third, probabilistic games can combine adversarial strategy
choices with stochastic transitions; local inequalities become nested
minimum/maximum expectation conditions, and qualitative recurrence requires
its own end-component analysis. The third extension is \emph{not} implemented
by merely combining the current parity and MDP parsers.

Another practical direction is policy compression. A finite table can be
synthesized into a small source-level decision tree using Leant/Djex or Forge's
existing witness facilities, then proved extensionally equal to the certified
table. This can improve the readability and executable size of extracted
controllers without forcing the game solver to search directly in a
higher-order term language. Compression is an optional post-pass: failure
leaves the original checked table usable.

\subsection{Assessment}

The proposed increase in capability is specific. Forge gains finite witnesses
for infinite adversarial behavior, exact relational correlations, quantified
matching of probabilistic actions, and expected termination in the presence of
cycles. These are proof shapes its existing universal-closure and pathwise-rank
workers do not supply by themselves.

The four implementations demonstrate the certificate computations and reject
important semantic corruptions. The formalization obligation is equally
specific: prove the local checker theorems and connect exact finite models to
original Lean expressions. The delivered evidence supports the former as a
mathematical design and an executed Python prototype, not yet as a Lean
implementation. A successful next contribution should make that missing bridge
small, checked, and reusable rather than enlarge the planner's unverified
surface.
